\section{Week 11}

\subsection{Caley's Theorem}

Last time we discussed \( G  \) acting on itself by left-multiplication which is a special case of \( G  \) acting on the cosets of a subgroup \( H  \) by left-multiplication; that is, \( g \cdot aH  = (ga)H \) for all \( g \in G  \) and all \( a H  \in G/ H   \). We will prove some things using the permutation representation of \( g  \) acting on the cosets of a subgroup \( H \). That is, 
\[  \varphi : G \to {S}_{G / H } \]
defined by \( \psi(g) = \sigma_{g} \) where \( {\sigma}_{g}(xH) = g \cdot xH =  (gx)H  \) is a homomorphism. 

\begin{theorem}[ ]
    Let \( G  \) and be a group and \( H \leq G  \). Let \( G  \) act on \( A = G / H  \) by left multiplication. Let \( {\pi}_{H} : G \to {S}_{G /H}  \) be the associated permutation representation afforded by this action. Then 
    \begin{enumerate}
        \item[(1)] \( G  \) acts transitively on \( A  \); that is, for any \( aH , b H \in A  \), there exist a \( g \in G  \) such that  
            \[  b H= g \cdot a H  \]
        \item[(2)] \( {G}_{1H} = H  \).
        \item[(3)] The kernel of the action (\( \text{ker}({\pi}_{H}) \)) is 
            \[  \bigcap_{  x \in G  }^{  }  x H x^{-1}  \]
            and the kernel of \( \pi_{H}  \) is the largest normal subgroup of \( G  \) contained in \( H  \).
    \end{enumerate}
\end{theorem}
\begin{proof}
\begin{enumerate}
    \item[(a)] To see that \( G  \) acts transitively on \( A  \), let \( a H , b H \in A  \). Then let \( g = b a^{-1}  \) and so 
        \begin{align*}
           g \cdot a H  &= (b a^{-1})\cdot aH  \\
                        &= (b a^{-1}) aH \\
                        &= b (a^{-1}a) H  \\
                        &= b H. 
        \end{align*}
    \item[(b)] Note that 
        \begin{align*}
            {G}_{1H} &= \{  g \in G : g \cdot 1 H  = 1 H  \}  \\
                     &= \{  g \in G : (g1)H = H  \}  \\
                     &= \{  g \in G : g H  = H  \} \\
                     &= \{  g \in G  : 1^{-1} g \in H  \}  \\
                     &= \{  g \in G  : g \in H  \}  \\
                     &= H. 
        \end{align*}
    \item[(c)] By definition, 
        \begin{align*}
            \text{ker}({\pi}_{H}) &= \{  g \in G : {\pi}_{H}(g) = {1}_{{S}_{G /H}} \}   \\
                                  &= \{  g \in G : g \cdot x H = x h \ \forall xH \in A  \} \\
                                  &= \{  g \in G  : g x H = x H \ \forall x \in G   \} \\ 
                                  &= \{  g \in G : x^{-1} g x \in H  \ \forall x \in G  \}  \\
                                  &= \{  g \in G : g \in x H x^{-1} \forall x \in G  \}  \\
                                  &= \bigcap_{ g \in G  }^{  } x H x^{-1}. 
    \end{align*}
\end{enumerate}
Observe that  \( \text{ker}({\pi}_{H}) \trianglelefteq G  \) and \( \text{ker}({\pi}_{H}) \leq H  \). Let \( N \trianglelefteq G  \) such that \( N \leq H  \) if and only if \( x N x^{-1} \subseteq  x H x^{-1} \). Then 
\[  N = x N x^{-1} \leq x H x^{-1} \ \ \forall  x \in G. \]
Hence, 
\[  N \leq \bigcap_{  x \in G  }^{  }  x H x^{-1} = \text{ker}({\pi}_{H}). \]

\end{proof}

\begin{theorem}[Caley's Theorem]
   Every group is isomorphic to a subgroup of some symmetric group. In particular, if \( | G  |  = n < \infty \), then \( G  \) is isomorphic to a subgroup of \( {S}_{n} \). 
\end{theorem}
\begin{proof}
Let \( H  = \{  1  \}  \) and by the first iso theorem, we have that 
\[  G / \text{ker}({\pi}_{H}) \cong {\pi}_{H}(G) \leq {S}_{G}. \]
Since \( \text{ker}({\pi}_{H}) \leq H  \), we have 
\[  \text{ker}({\pi}_{H}) \leq \{  1 \} \implies  \text{ker}({\pi}_{H}) = 1  .  \]
Hence, 
\[   G \cong  G / \text{ker}({\pi}_{H}) \cong {\pi}_{H}(G).  \]
where \( {\pi}_{H}(G) \leq {S}_{G} \).
\end{proof}

The permutation representation of this action when \( H = \{  1  \}  \) is called the left regular representation of \( G  \). 

\begin{corollary}[Classifying Groups]
   If \( G  \) is a finite group of order \( n  \in \N  \) and \( p  \) is the smallest prime dividing \( n  \), then any subgroup of index \( p  \) is normal in \( G  \). 
\end{corollary}
\begin{proof}
Suppose \( H \leq G  \) and \( | G : H  |  = p  \). Let \( {\pi}_{H}  \) be the permutation representation afforded by \( G  \) acting on the left cosets of \( H  \) by left-multiplication. Let \( K = \text{ker}({\pi}_{H}) \) and \( | H : K  |  = k  \). Notice that \( k \mid | H  |  \) (by Lagrange's Theorem) and so \( k \mid | G  |  \). Then by our tower laws, we see that
\begin{align*}
    | G : K  | &=| G : H  | | H : K  | = p k.  
\end{align*}
Since \( H  \) has \( p  \) cosets in \( G  \), 
\[  \pi_{H} : G \to {S}_{p}. \]
By the first isomorphism theorem, 
\[  G / K \cong {\pi}_{H}(G) \leq {S}_{p}. \]
This implies that \( pk = | G :  K | = \frac{  | G  |  }{ | K  |  }    \). Hence, 
\[  pk  = | {\pi}_{H}(G) |  \]
and so 
\[ p k \mid | {S}_{p} | \implies pk \mid p^{!} \implies k \mid (p-1)!. \]
However, since \( k \mid | G  |  \) any prime factor of \( k  \).  However, since \( k \mid | G  |  \) any prime factors of \( k  \) would be larger than or equal to \( p  \), so \( k = 1  \) minimality of \( p  \). Indeed, suppose otherwise. Then \( q \mid k   \) implies \( q \mid | G  |  \) and \( q \mid (p-1)! \). But this tells us that since \( | H : K  |  = k  \), we have \( | H : K  |  = 1  \) and so \( H = K  \). Hence, \( H = K \trianglelefteq G  \).
\end{proof}

\begin{eg}
    Given a group \( G  \) of order \( 15  \) by Cauchy's Theorem, we know that \( G  \) contains a subgroup of order \( 5  \), and by the previous corollary, we know that subgroup is normal. Since its index must be \( 3  \) which is the smallest prime dividing \( 15 \).
\end{eg}

\subsection{Groups Acting on Themselves}

Groups acting on themselves by conjugation (class equation). Let \( G  \) be a group acting on itself by conjugation. Define \( g \cdot a =  g a g^{-1} \) for all \( g \in G  \) and all \( a \in G  \). We proved that this is a group action.

\begin{definition}[ ]
    We say \( a,b \in G  \) are conjugate in \( G  \) if there exists a \( g \in G  \) such that 
    \[  b = g a g^{-1}; \]
    that is, \( a,b  \) are conjugate in \( G  \) if and only if \( a,b \in {\mathcal{O}}_{x}  \) for some \( x \in G  \).
Recall that 
\[  {\mathcal{O}}_{x} = \{  g \cdot x : g \in G  \} = \{ g x g^{-1} : g\in  G    \}. \]
The orbits of \( G  \) acting on itself by conjugation are called the \textbf{conjugacy classes} of \( G  \).
\end{definition}

\begin{eg}
   \begin{itemize}
       \item If \( G  \) is abelian, then \( G  \) acting on itself by conjugation is the  \textbf{trivial action}; that is, \( g \cdot a = g a g^{-1} = a  \) for all \( a \in G  \)  and \( g \in G  \).  
   \end{itemize} 
\end{eg}
This tells us that the conjugacy class of \( a  \) in \( G \) is
\begin{align*}
    \{  g a g^{-1 } : g \in G  \} &= \{ a :  g \in G  \}  \\
                                  &= \{  a  \}.
\end{align*}
More generally, for \( G  \) acting on itself where \( G   \) is not assumed to be abelian, the conjugacy class of some \( g \in G  \) is just \( \{  g  \}  \) if and only if \( g  \in Z(G)  \).
\begin{eg}
    If \( | G  |   > 1  \), then \( G  \) does not act transitively on itself by conjugation. (If transitive, there is only one orbit) The conjugacy class of \( {1}_{G} \) is just \( \{  {1}_{G} \}  \). That is, 
    \[  \{  g 1 g^{-1} : g \in G  \} = \{  1 : g \in G  \}  = \{ 1  \}. \]
\end{eg}

Note that also the action of conjugation can be generalized to \( G  \) acting on \( \mathcal{P}(G) \) by
\[  g \cdot S = g S g^{-1} \ \forall g \in G ,\ S \in \mathcal{P}(G) \]
where \( \mathcal{P} \) is the power set of \( G \). 

\begin{definition}[Conjugate]
    Two subsets \( S  \) and \( T  \) of \( G  \) are called \textbf{conjugate} in \( G  \) if there exists a \( g \in G  \) such that \( S = g T g^{-1} \). 
\end{definition}

\begin{lemma}[Result Needed to Prove Class Equation (Prop 6)]
The number of conjugates of a subset \( S  \) in a group \( G  \) is \( | G : {N}_{G}(S) |  \). In particular, the number of conjugates of an element \( s \in G  \) is \( | G : {C}_{G}(s) |  \).    
\end{lemma}
\begin{proof}
Let \( S \subseteq  G  \). Using Orbit-stabilizer theorem, we know that the conjugacy class of \( S  \) equals \( | G : {G}_{s} |  \). For this particular action,  
\begin{align*}
    {G}_{s} &= \{  g \in G : g \cdot S = S  \}  \\
            &= \{  g \in G : g S g^{-1} = S  \} \\
            &= {N}_{G}(S).
\end{align*}
The second statement follows from 
\begin{align*}
    {N}_{G}(\{ s \} ) &= \{  g \in G : g \{ s  \}  g^{-1} = \{ s \}  \}  \\
                      &= \{  g \in G : g s g^{-1} = s  \}  \\
                      &= {C}_{G}(\{ s \}).
\end{align*}
That is, \( {N}_{G}(s) = {C}_{G}(s) \).
\end{proof}

\begin{theorem}[The Class Equation]
   Let \( G  \) be a finite group and let \( {g}_{1}, {g}_{2}, \dots , {g}_{r} \) be representatives of distinct distinct conjugacy classes of \( G  \) not contained in \( Z(G) \). Then  
   \[  | G  |  = | Z(G) |  + \sum_{ i=1  }^{ r } | G : {C}_{G}({g}_{i}) |. \]
\end{theorem}
\begin{proof}
We claim that \( z \in G \) has conjugacy class \( \{  z  \}  \) if and only if \( z \in Z(G) \). You can show this in your free time. Let 
\[  Z(G) = \{  1, {z}_{2}, {z}_{3}, \dots, {z}_{m} \}.  \]
Let \( {k}_{1}, {k}_{2}, \dots, {k}_{r} \) be the conjugacy classes of \( G  \) not contained in \( Z(G) \) and let \( {g}_{i} \) be a representative of \( {k}_{i} \) (element of) for each \( i  \). Then the full set of conjugacy classes is \[ \{ 1 \} , \{ {z}_{2} \}, \dots, \{ {z}_{m} \}, {k}_{1}, {k}_{2}, \dots, {k}_{r}. \]
Since these partition \( G  \), we have 
\begin{align*}
    | G  |  &= \sum_{ i=1  }^{ m } 1 + \sum_{ i=1  }^{ r } | {K}_{i} |  \\
            &= | Z(G) |  + \sum_{ i=1  }^{ r  } | G : {C}_{G}({g}_{i}) | 
\end{align*}
by the previous proposition.
\end{proof}

\begin{eg}
   \begin{enumerate}
       \item[(1)] If \( G  \) is abelian of order \( n  \), then the class equation tells us nothing; that, is
           \[  | G  |  = | Z(G) |. \]
        \item[(2)] Calculate the class equation of \( {Q}_{8} \). Note that in any group \( \langle  x  \rangle \leq {C}_{G}(x) \). We know that \( i \notin Z(G) \) (to see this compute \( i \) and \( j  \) together in different order). We also know that \( | \langle  i \rangle |  = 4  \) which implies that \( | {Q}_{8} : \langle  i  \rangle |  = 2  \). Also, we know that \( i \notin Z({Q}_{8}) \). This implies that the \( {C}_{{Q}_{8}}(i) < {Q}_{8}  \) (proper subgroup of \( {Q}_{8} \)). (Recall that \( {C}_{G}(x)  = G \iff x \in Z(G) \)). So, we have 
            \begin{align*}
                2 = | {Q}_{8} : \langle  i \rangle | &= | {Q}_{8} : {C}_{{Q}_{8}}(i)  |  | {C}_{{Q}_{8}}(i) : \langle i   \rangle    |
            \end{align*}
            If \( | {Q}_{8} : {C}_{{Q}_{8}} |  = 1  \), then \( {Q}_{8} = {C}_{{Q}_{8}}(i) \) is a contradiction because we know that \( i \notin Z(G) \). Thus, \( | {C}_{{Q}_{8}}(i) : \langle  i   \rangle |  = 1  \) and \( | {Q}_{8} : {C}_{{Q}_{8}}(i) | = 2  \). Thus, see that \( {C}_{{Q}_{8}} = \langle  i \rangle \). Now, it follows that the conjugacy class of \( i  \) contains an index of \( 2  \) elements, namely, \( i  \) and \( - i  \). Similarly, we can perform the same process to find all the conjugacy classes of \( {Q}_{8} \), 
            \[  \{  1  \} , \{  - 1  \} , \{  \pm i  \} , \{  \pm j  \} , \{  \pm k  \}. \]
            Thus, the class equation of \( {Q}_{8} \) is
            \begin{align*}
                | \{ 1  \}  |  + | \{  - 1  \}  |  + | \{  \pm i  \}   |  + | \{  \pm j  \}  |  + | \{  \pm k \}  |  =  2 + 2  +2 + 2 &= 8  \\
            \end{align*}
   \end{enumerate} 
\end{eg}

\begin{theorem}[ ]
    If \( p  \) is a prime and \( P  \) is a group of order \( p^{\alpha} \) with \( \alpha \geq 1  \). Then \( P  \) has a non-trivial center; that is, \( Z(P) \neq \{ {1 }_{P} \}   \) or \( | Z(P) |  > 1  \). 
\end{theorem}

\begin{proof}
    By the class equation, 
    \[ | P  |  = | Z(P)    |  + \sum_{ i =1  }^{ r  } | P : {C}_{P}({g}_{i}) |  \]
    where \( {g}_{1}, {g}_{2}, \dots, {g}_{r} \) are representatives of the distinct non-central conjugacy classes. By definition, \( {C}_{P}({g}_{i}) \neq P  \). (Note that \( {C}_{P}({g}_{i}) = P  \) if and only if \( {g}_{i} \in Z(G) \)) for \( i = 1,2, \dots, r  \). So, \( p  \) divides the \[ | P : {C}_{P}({g}_{i}) | = \frac{ | P |  }{  | {C}_{P}({g}_{i}) |  } = \frac{ p^{\alpha} }{  | {C}_{p}({g}_{i}) |  } = p^{\beta}    \]
    where \( 0  < \beta \leq \alpha \). Since \( p  \) also divides \( |  P  |   \) it follows that 
    \[  p \mid \Big(  | P  |  - \sum_{ i=1  }^{ r } | P : {C}_{P}({g}_{i}) |  \Big) \]
    equivalently, we have \( p \mid | Z(P) |  \). Hence, it follows that \( Z(G)  \) is non-trivial.
\end{proof}

\begin{corollary}
    If \( | P  | = p^{2} \) where \( p  \) is a prime, then \( P  \) is abelian. More precisely, \( P \cong {Z}_{p^{2}}  \) or \( P \cong {Z}_{p} \times {Z}_{p} \). 
\end{corollary}
\begin{proof}
Since \( Z(P) \neq 1  \) from the previous theorem, it follows that \( | Z(P) |  \in \{  p, p^{2} \}  \) and so we have 
\[  |  P / Z(P) | \in \{  1, p  \}.   \]
In either case, \( P   \) is abelian (see result from one of homeworks; that is, \( G / Z(G) \) implies that  G  is abelian). By the above, \( P  \) is abelian. If \( P  \) has an element of order \( p^{2} \). If \( P  \) has an element of order \( p^{2} \), \( P  \) is cyclic and isomorphic to \( {Z}_{p^{2}} \). If \( P \) contains no elements of order \( p^{2} \), then every non-identity element of \( P \) has order \( p  \). Let \( x \in P  \) such that \( x \neq 1  \) and \( y \in P - \langle  x  \rangle \) (\( y \neq 1  \) since \( 1 \in \langle  x \rangle \)). Notice that \( x  \) and \( y  \) both have order \( p  \). Since \( | \langle x , y \rangle |  > | \langle x   \rangle | = p  \), we have \( | \langle x , y \rangle |  = p^{2} \); that is, \( \langle x , y \rangle = P  \). Note since \( P  \) is abelian, \( \langle x  ,  y  \rangle = \{  x^{a} y^{b } : a,b \in \Z \}  \). Also, since \( | x  |  = | y  |  = p  \), \( \langle x  \rangle \times \langle y  \rangle  \cong {Z}_{p} \times {Z}_{p} \). Defining a map from \( {Z}_{p} \times {Z}_{p} \) to \( P  \) by \( (x^{a}, y^{b}) \mapsto x^{a} y^{b} \) gives an isomorphism. Hence, \( {Z}_{p} \times {Z}_{p} \cong P  \).
\end{proof}




